\chapter{Proiectare de detaliu și implementare}\label{ch:proiectare}
\pagestyle{fancy}

Acest capitol prezintă proiectarea detaliată și implementarea în VHDL a
sistemului de control FPGA pentru mâna robotică. Fiecare modul hardware este
descris din perspectiva funcționalității, a interfeței de porturi, a mașinii de
stări interne și a deciziilor de proiectare adoptate. Codul sursă complet este
inclus în anexe; în acest capitol se prezintă doar fragmente relevante,
însoțite de explicații funcționale.

%==========================================================================
\section{Schema generală a sistemului}
\label{sec:schema-generala}

Sistemul este compus din 12 module VHDL interconectate ierarhic prin intermediul
modulului de nivel superior \texttt{basys3\_top}. Figura~\ref{fig:block-diagram}
prezintă diagrama bloc a întregului sistem, evidențiind fluxul de date de la
intrarea UART până la ieșirile PWM.

\begin{figure}[H]
\centering
\fbox{\parbox{0.95\textwidth}{\centering\vspace{6pt}
\small
\texttt{uart\_rx} $\rightarrow$ \texttt{frame\_rx} $\rightarrow$
\texttt{mailbox} $\rightarrow$ \texttt{source\_mux} $\rightarrow$
\texttt{kalman\_engine} $\rightarrow$ \texttt{pwm\_gen} (8 instanțe)\\[4pt]
\texttt{system\_tick\_gen} --- ceas 20\,ms pentru toate modulele\\[4pt]
\texttt{demo\_rom} $\rightarrow$ \texttt{source\_mux} (intrare alternativă)\\[4pt]
\texttt{spi\_master} $\leftarrow$ \texttt{ad7124\_ctrl} $\rightarrow$
\texttt{kalman\_engine} (feedback ADC)\\[4pt]
\texttt{frame\_tx} $\rightarrow$ \texttt{uart\_tx} (răspuns către PC)
\vspace{6pt}}}
\caption{Diagrama bloc a sistemului --- fluxul de interconectare al modulelor VHDL}
\label{fig:block-diagram}
\end{figure}

Cele 12 module sunt enumerate în Tabelul~\ref{tab:module-list}, împreună cu
rolul funcțional al fiecăruia.

\begin{table}[H]
\centering
\caption{Lista modulelor VHDL și rolul lor în sistem}
\label{tab:module-list}
\begin{tabular}{@{}lp{9cm}@{}}
\toprule
\textbf{Modul} & \textbf{Rol funcțional} \\
\midrule
\texttt{system\_tick\_gen} & Generează pulsul periodic de 20\,ms \\
\texttt{uart\_rx}          & Recepția serială UART la 115\,200\,baud \\
\texttt{uart\_tx}          & Transmisia serială UART la 115\,200\,baud \\
\texttt{frame\_rx}         & Parser de cadre de 10\,octeți cu verificare checksum \\
\texttt{frame\_tx}         & Serializare cadru de răspuns (direcție 0xFE) \\
\texttt{mailbox}           & Buffer dublu (ping-pong) pentru decuplare temporală \\
\texttt{source\_mux}       & Multiplexor combinațional: UART vs. demo \\
\texttt{demo\_rom}         & ROM cu 21 de cadre-cheie pentru mod demonstrativ \\
\texttt{spi\_master}       & Master SPI mod~3, 5\,MHz \\
\texttt{ad7124\_ctrl}      & Controler pentru ADC-ul AD7124-8 \\
\texttt{kalman\_engine}    & Filtru Kalman cu 2 stări, Q16.16, 8 canale \\
\texttt{pwm\_gen}          & Generator PWM pentru un servo (8 instanțe) \\
\bottomrule
\end{tabular}
\end{table}

Interconexiunile dintre module respectă principiul fluxului de date
unidirecțional: datele curg de la recepția UART către generarea PWM, cu o
ramură de feedback de la ADC prin filtrul Kalman. Semnalul de ceas de 20\,ms
(\texttt{tick\_20ms}) sincronizează toate modulele, asigurând că actualizarea
servomotoarelor se face cu o perioadă fixă, independentă de rata de sosire a
cadrelor UART.

%==========================================================================
\section{Modulul \texttt{system\_tick\_gen}}
\label{sec:system-tick-gen}

Modulul \texttt{system\_tick\_gen} generează un puls de sincronizare cu
perioada exactă de 20\,ms, necesară pentru ciclul standard al servomotoarelor
(50\,Hz). Implementarea se bazează pe un numărător pe 21 de biți care
numără de la 0 la 1\,999\,999, corespunzând celor 2\,000\,000 de cicluri
de ceas la 100\,MHz necesare pentru 20\,ms.

\subsection{Interfața de porturi}

\begin{table}[H]
\centering
\caption{Porturile modulului \texttt{system\_tick\_gen}}
\begin{tabular}{@{}llp{7.5cm}@{}}
\toprule
\textbf{Port} & \textbf{Direcție} & \textbf{Descriere} \\
\midrule
\texttt{clk}       & intrare  & Ceasul de sistem, 100\,MHz \\
\texttt{rst}       & intrare  & Reset sincron \\
\texttt{tick\_20ms} & ieșire   & Puls de un ciclu de ceas la fiecare 20\,ms \\
\texttt{tick\_cnt}  & ieșire   & Valoarea curentă a numărătorului (21\,biți) \\
\bottomrule
\end{tabular}
\end{table}

\subsection{Funcționare}

Numărătorul \texttt{counter} este de tip \texttt{unsigned(20 downto 0)} și se
incrementează la fiecare front crescător al ceasului. Când atinge valoarea
constantei \texttt{PERIOD} = 1\,999\,999, se resetează la zero. Ieșirea
\texttt{tick\_20ms} este '1' exclusiv în ciclul de ceas în care numărătorul
egalează \texttt{PERIOD}, generând astfel un puls de exactmente un ciclu de
ceas (10\,ns) la fiecare 20\,ms.

Ieșirea \texttt{tick\_cnt} furnizează valoarea curentă a numărătorului,
utilizată de modulul \texttt{pwm\_gen} pentru compararea cu pragul PWM.
Deoarece ambele module partajează același numărător de bază, semnalele PWM
sunt perfect sincronizate cu ciclul de 20\,ms.

Calcul: $100\,\text{MHz} \times 20\,\text{ms} = 100 \times 10^6 \times 20
\times 10^{-3} = 2\,000\,000$ cicluri. Numărătorul numără $0, 1, \ldots,
1\,999\,999$, deci necesită $\lceil \log_2(2\,000\,000) \rceil = 21$ biți.

%==========================================================================
\section{Modulul \texttt{uart\_rx} --- Recepția UART}
\label{sec:uart-rx}

Modulul \texttt{uart\_rx} implementează un receptor UART complet la
115\,200\,baud, cu protecție contra metastabilității și eșantionare la mijlocul
bitului. Configurația este 8N1 (8 biți de date, fără paritate, 1 bit de stop).

\subsection{Calculul parametrilor de baud}

Rata de ceas a sistemului este de 100\,MHz, iar rata de transfer (baud rate)
este de 115\,200\,baud. Numărul de cicluri de ceas pe bit este:

\begin{equation}
\text{CLKS\_PER\_BIT} = \frac{100\,000\,000}{115\,200} - 1 = 868 - 1 = 867
\label{eq:clks-per-bit}
\end{equation}

Scăderea cu 1 se datorează faptului că numărătorul pornește de la 0. Pentru
eșantionarea la mijlocul bitului, se utilizează constanta:

\begin{equation}
\text{HALF\_BIT} = \left\lfloor \frac{868}{2} \right\rfloor - 1 = 433
\label{eq:half-bit}
\end{equation}

Eșantionarea la mijlocul bitului este esențială deoarece semnalul serial
prezintă tranziții la granițele biților, iar valorile sunt stabile doar în
zona centrală. Prin așteptarea a 433 de cicluri după detectarea frontului
descrescător al bitului de start, receptorul se poziționează în centrul
fiecărui bit, maximizând marja de toleranță la variațiile de frecvență
între emițător și receptor.

\subsection{Protecția contra metastabilității}
\label{subsec:metastability}

Semnalul \texttt{rx\_serial} provine dintr-un domeniu de ceas extern (linia
UART a PC-ului) și este asincron față de ceasul FPGA de 100\,MHz. Când un
semnal asincron este eșantionat de un flip-flop, există o probabilitate
nenulă ca acesta să fie capturat exact în momentul tranziției, rezultând o
stare metastabilă --- o tensiune intermediară între '0' și '1' care nu se
rezolvă în timpul de setup al flip-flop-ului. Această situație poate provoca
comportament nedeterminist în logica din aval.

Soluția standard constă în utilizarea a două flip-flop-uri în cascadă,
cunoscute sub numele de \textit{synchronizer} sau \textit{double-FF}:

\begin{lstlisting}[caption={Sincronizator cu dublu flip-flop pentru protecția contra metastabilității},label={lst:metastab}]
signal rx_sync1 : std_logic := '1';
signal rx_sync2 : std_logic := '1';

process(clk)
begin
    if rising_edge(clk) then
        rx_sync1 <= rx_serial;
        rx_sync2 <= rx_sync1;
    end if;
end process;
\end{lstlisting}

Primul flip-flop (\texttt{rx\_sync1}) poate intra în stare metastabilă, dar
are un ciclu de ceas complet (10\,ns la 100\,MHz) pentru a se rezolva înainte
ca al doilea flip-flop (\texttt{rx\_sync2}) să-l eșantioneze. Probabilitatea
ca metastabilitatea să se propage prin ambele flip-flop-uri este neglijabilă
--- de ordinul $10^{-20}$ erori pe secundă pentru flip-flop-urile din
tehnologia Artix-7.

Valorile inițiale sunt '1' deoarece linia UART este inactivă la nivel logic
înalt (\textit{idle high}).

\subsection{Mașina de stări}

Receptorul UART utilizează o mașină de stări cu 4 stări:

\begin{enumerate}
	\item \textbf{S\_IDLE} --- Stare de așteptare. Monitorizează linia
	      \texttt{rx\_sync2} pentru detectarea frontului descrescător care
	      marchează începutul bitului de start.
	\item \textbf{S\_START} --- Verificarea bitului de start. Numărătorul de
	      ceas (\texttt{clk\_count}) avansează până la \texttt{HALF\_BIT} (433),
	      moment în care se verifică dacă linia este încă '0'. Dacă da,
	      tranziția a fost legitimă și se trece la recepția datelor; dacă nu,
	      a fost un impuls parazit (\textit{glitch}) și se revine la S\_IDLE.
	\item \textbf{S\_DATA} --- Recepția celor 8 biți de date. La fiecare
	      \texttt{CLKS\_PER\_BIT} (867) cicluri, valoarea de pe linie este
	      memorată în registrul de deplasare la poziția indicată de
	      \texttt{bit\_idx}. Biții sunt recepționați LSB first, conform
	      standardului UART. După bitul 7, se trece la S\_STOP.
	\item \textbf{S\_STOP} --- Așteptarea bitului de stop. Se așteaptă
	      \texttt{CLKS\_PER\_BIT} cicluri, timp în care linia trebuie să fie
	      la nivel logic '1'. La finalul perioadei, \texttt{rx\_valid} este
	      activat pentru un ciclu de ceas, iar \texttt{rx\_data} conține
	      octetul recepționat complet.
\end{enumerate}

Durata totală de recepție a unui octet este: 1~bit~start + 8~biți~date +
1~bit~stop = 10~biți $\times$ 868~cicluri = 8\,680~cicluri $\approx$
86,8\,$\mu$s, corespunzând exact unui caracter la 115\,200~baud.

%==========================================================================
\section{Modulul \texttt{uart\_tx} --- Transmisia UART}
\label{sec:uart-tx}

Modulul \texttt{uart\_tx} este echivalentul simetric al receptorului, realizând
transmisia serială la 115\,200~baud. La activarea semnalului \texttt{tx\_start},
octetul prezent pe portul \texttt{tx\_data} este salvat în registrul intern
\texttt{data\_reg}, iar transmisia începe.

\subsection{Mașina de stări}

Transmițătorul utilizează aceleași 4 stări ca receptorul:

\begin{enumerate}
	\item \textbf{S\_IDLE} --- Linia de transmisie este menținută la '1' (idle).
	      La detectarea semnalului \texttt{tx\_start}, octetul este salvat și
	      transmisia începe.
	\item \textbf{S\_START} --- Linia este forțată la '0' pentru durata unui bit
	      complet (868~cicluri de ceas), semnalizând receptorului începutul unui
	      caracter.
	\item \textbf{S\_DATA} --- Cei 8~biți de date sunt transmiși secvențial,
	      câte un bit la fiecare 868~cicluri, începând cu bitul cel mai puțin
	      semnificativ (LSB first). Semnalul \texttt{tx\_serial} ia valoarea
	      bitului curent din \texttt{data\_reg(bit\_idx)}.
	\item \textbf{S\_STOP} --- Linia revine la '1' pentru durata unui bit,
	      marcând bitul de stop. La final, mașina de stări revine la S\_IDLE.
\end{enumerate}

Semnalul \texttt{tx\_busy} indică faptul că transmițătorul este ocupat
(orice stare diferită de S\_IDLE). Modulele din amonte (în particular
\texttt{frame\_tx}) verifică acest semnal înainte de a iniția o nouă
transmisie, prevenind suprapunerea caracterelor.

Spre deosebire de receptor, transmițătorul nu necesită protecție contra
metastabilității, deoarece atât datele de intrare cât și semnalul de ieșire
se află în același domeniu de ceas de 100\,MHz.

%==========================================================================
\section{Modulul \texttt{frame\_rx} --- Parserul de cadre}
\label{sec:frame-rx}

Modulul \texttt{frame\_rx} asamblează cadre de 10~octeți din octeții
individuali recepționați de \texttt{uart\_rx}. Protocolul de cadre definit
în Capitolul~\ref{ch:analysis} specifică formatul:

\begin{center}
\begin{tabular}{|c|c|c|}
\hline
Octet 0 & Octeții 1--8 & Octet 9 \\
\hline
Direcție (0xFF/0xFE) & 8 unghiuri servo (0--180) & Checksum XOR \\
\hline
\end{tabular}
\end{center}

\subsection{Mașina de stări}

Parserul de cadre utilizează 3~stări:

\begin{enumerate}
	\item \textbf{S\_WAIT\_DIR} --- Așteaptă un octet de direcție valid
	      (0xFF pentru scriere unghiuri, 0xFE pentru citire). Orice alt octet
	      este ignorat, oferind o resincronizare naturală a protocolului.
	      La primirea unui octet valid, acesta este salvat în \texttt{dir\_reg}
	      și este inițializat acumulatorul XOR (\texttt{xor\_acc}) cu valoarea
	      direcției.

	\item \textbf{S\_RECV\_DATA} --- Colectează 8~octeți de date (unghiurile
	      servomotoarelor). Fiecare octet recepționat este stocat în registrul
	      \texttt{angles\_reg} la poziția corespunzătoare și este acumulat în
	      \texttt{xor\_acc} prin operația XOR. Numărătorul \texttt{byte\_count}
	      urmărește câți octeți au fost primiți. După al 8-lea octet, se trece
	      la S\_RECV\_CHKSUM.

	\item \textbf{S\_RECV\_CHKSUM} --- Recepționează octetul de checksum.
	      Verificarea se face prin calculul
	      $\texttt{xor\_acc} \oplus \texttt{rx\_data}$: dacă rezultatul este
	      0x00, cadrul este valid și semnalul \texttt{frame\_valid} este activat
	      pentru un ciclu de ceas; în caz contrar, \texttt{checksum\_err} este
	      activat. În ambele cazuri, mașina revine la S\_WAIT\_DIR.
\end{enumerate}

\subsection{Mecanismul de timeout}

Dacă un cadru nu este completat în 1\,ms (100\,000~cicluri de ceas), modulul
revine automat la starea S\_WAIT\_DIR. Acest mecanism previne blocarea în
cazul pierderii de octeți pe linia serială. Numărătorul de timeout
(\texttt{timeout\_cnt}) pe 17~biți se incrementează în fiecare ciclu de ceas
când starea nu este S\_WAIT\_DIR și se resetează la zero la fiecare octet
recepționat valid. Valoarea pragului este
\texttt{TIMEOUT\_MAX} = 99\,999 ($\approx$ 1\,ms la 100\,MHz).

\subsection{Principiul verificării XOR}

Emițătorul calculează checksum-ul ca $\bigoplus_{i=0}^{8} b_i$ (XOR-ul tuturor
octeților de la 0 la 8) și îl plasează în octetul 9. Receptorul acumulează
XOR-ul tuturor celor 10~octeți: dacă niciun bit nu a fost alterat în
transmisie, rezultatul este $\bigoplus_{i=0}^{9} b_i = \bigoplus_{i=0}^{8} b_i
\oplus b_9 = x \oplus x = 0$. Orice bit inversat în tranzit produce un
rezultat nenul, iar cadrul este rejectat.

%==========================================================================
\section{Modulul \texttt{frame\_tx} --- Transmisia cadrelor de răspuns}
\label{sec:frame-tx}

Modulul \texttt{frame\_tx} serializează un cadru de răspuns atunci când
calculatorul gazdă trimite un cadru cu direcția 0xFE (cerere de citire).
Răspunsul conține cele 8~unghiuri filtrate curent, precedate de octetul de
direcție 0xFE și urmate de checksum-ul XOR.

\subsection{Mașina de stări}

Transmisia se desfășoară în 7~stări, grupate în 3~faze:

\begin{enumerate}
	\item \textbf{S\_IDLE} --- Așteaptă activarea semnalului \texttt{send\_start}.
	      La activare, unghiurile sunt salvate în \texttt{angles\_reg}.

	\item \textbf{S\_SEND\_DIR / S\_WAIT\_DIR} --- Transmite octetul de
	      direcție 0xFE prin \texttt{uart\_tx}. Inițializează acumulatorul XOR
	      cu 0xFE. Așteaptă ca \texttt{tx\_busy} să devină '0' înainte de a
	      continua.

	\item \textbf{S\_SEND\_DATA / S\_WAIT\_DATA} --- Transmite secvențial cei
	      8~octeți de unghiuri, acumulând checksum-ul XOR. Indexul
	      \texttt{byte\_idx} contorizează octeții transmiși, de la 0 la 7.

	\item \textbf{S\_SEND\_CHKSUM / S\_WAIT\_CHKSUM} --- Transmite octetul de
	      checksum (\texttt{xor\_acc}) și revine la S\_IDLE.
\end{enumerate}

Semnalul \texttt{frame\_busy} indică faptul că transmisia unui cadru este
în curs. Modulul de nivel superior verifică acest semnal pentru a preveni
inițierea simultană a mai multor cadre de răspuns.

Între fiecare octet transmis, modulul verifică semnalul \texttt{tx\_busy}
de la \texttt{uart\_tx} pentru a se asigura că transmițătorul serial a
terminat de trimis octetul precedent. La 115\,200~baud, transmisia unui
octet complet (10~biți) durează aproximativ 86,8\,$\mu$s, iar un cadru
complet de 10~octeți durează aproximativ 868\,$\mu$s.

%==========================================================================
\section{Modulul \texttt{mailbox} --- Buffer dublu}
\label{sec:mailbox}

Modulul \texttt{mailbox} implementează un mecanism de decuplare temporală
între rata de sosire a cadrelor UART (variabilă, dependentă de procesarea
MediaPipe pe calculatorul gazdă) și ciclul fix de 20\,ms al servomotoarelor.

\subsection{Principiul ping-pong}

Modulul conține două registre de 64~de biți:

\begin{itemize}
	\item \texttt{buf\_write} --- registrul de scriere, actualizat oricând
	      sosește un cadru UART valid.
	\item \texttt{buf\_read} --- registrul de citire, actualizat exclusiv pe
	      frontul semnalului \texttt{tick\_20ms}.
\end{itemize}

Un semnal de stare \texttt{new\_pending} indică dacă un cadru nou a fost
depus de la ultimul tick de 20\,ms. La fiecare tick, dacă
\texttt{new\_pending} este activ, conținutul \texttt{buf\_write} este copiat
în \texttt{buf\_read}, iar \texttt{new\_pending} este resetat. Ieșirea
\texttt{rd\_new} pulsează pentru un ciclu de ceas la fiecare actualizare.

Acest mecanism oferă două garanții importante:

\begin{enumerate}
	\item \textbf{Menținerea poziției:} Dacă niciun cadru nou nu sosește
	      într-un ciclu de 20\,ms, servomotoarele mențin ultima poziție
	      comandată. Mâna robotică rămâne stabilă chiar dacă comunicarea
	      UART este temporar întreruptă.
	\item \textbf{Rata de eșantionare:} Dacă mai multe cadre UART sosesc
	      în același interval de 20\,ms, doar ultimul este utilizat. Aceasta
	      este corectă deoarece doar cel mai recent unghi comandat este
	      relevant pentru ciclul curent.
\end{enumerate}

Condiția de scriere a mailbox-ului în modulul de nivel superior este:
\texttt{mb\_wr\_valid <= frame\_valid and (not frame\_chksum\_err)}, asigurând
că doar cadrele cu checksum valid sunt acceptate.

%==========================================================================
\section{Modulul \texttt{source\_mux} --- Multiplexor de surse}
\label{sec:source-mux}

Modulul \texttt{source\_mux} este un multiplexor pur combinațional care
selectează sursa unghiurilor țintă pentru filtrul Kalman:

\begin{lstlisting}[caption={Implementarea multiplexorului de surse},label={lst:source-mux}]
target_angles <= demo_angles when demo_mode = '1'
                 else mailbox_angles;
\end{lstlisting}

Când comutatorul \texttt{sw0} de pe placa Basys3 este în poziția ON
(\texttt{demo\_mode = '1'}), unghiurile provin din modulul \texttt{demo\_rom}.
În caz contrar, unghiurile provin din \texttt{mailbox}, deci din cadrele
UART recepționate de la calculatorul gazdă.

Logica este pur combinațională, fără registre, producând o latență de
propagare de ordinul nanosecundelor. Această decizie de proiectare este
adecvată deoarece semnalul \texttt{sw0} se modifică rar (doar la acțiunea
manuală a utilizatorului) și nu necesită filtrare de debounce --- un
comutator basculant nu prezintă fenomenul de bouncing, spre deosebire de
un buton.

%==========================================================================
\section{Modulul \texttt{demo\_rom} --- Secvențe demonstrative}
\label{sec:demo-rom}

Modulul \texttt{demo\_rom} conține o memorie ROM cu 21~de cadre-cheie
(\textit{keyframes}) care definesc secvențe de mișcare predefinite ale mâinii
robotice. Fiecare cadru-cheie conține 8~unghiuri (câte un octet per servo,
totalul de 64~biți) și o durată de menținere exprimată în tick-uri de 20\,ms.

\subsection{Formatul cadru-cheie}

Fiecare intrare în ROM este definită ca un \texttt{record} VHDL:

\begin{lstlisting}[caption={Tipul record pentru un cadru-cheie din ROM},label={lst:keyframe}]
type keyframe_t is record
    angles : std_logic_vector(63 downto 0);
    hold   : unsigned(7 downto 0);
end record;
\end{lstlisting}

Câmpul \texttt{hold} specifică numărul de tick-uri de 20\,ms cât se mențin
unghiurile respective. De exemplu, \texttt{hold = 50} corespunde unei
durate de $50 \times 20\,\text{ms} = 1\,\text{secundă}$.

\subsection{Secvențele de mișcare}

ROM-ul conține următoarele secvențe, ordonate astfel:

\begin{table}[H]
\centering
\caption{Secvențele demonstrative stocate în ROM}
\label{tab:demo-sequences}
\begin{tabular}{@{}clc@{}}
\toprule
\textbf{Indici} & \textbf{Secvența} & \textbf{Durată} \\
\midrule
0--3   & Deschidere/închidere pumn (2$\times$) & 4\,s \\
4--8   & Fluturare (rotire încheietură stânga/dreapta) & 2,5\,s \\
9--13  & Flexii individuale degete & 3\,s \\
14--15 & Gest \textit{thumbs up} & 2,5\,s \\
16--17 & Gest \textit{pointing} & 2,5\,s \\
18--20 & Flexia/extensia cotului & 3\,s \\
\bottomrule
\end{tabular}
\end{table}

Constanta \texttt{LAST\_KF = 20} definește ultimul cadru-cheie valid.
La depășirea acestui indice, numărătorul de cadre-cheie \texttt{kf\_index}
se resetează la zero, creând un ciclu continuu de demonstrație.

\subsection{Logica de avansare}

Un numărător \texttt{hold\_counter} se incrementează la fiecare tick de
20\,ms. Când atinge valoarea \texttt{hold} a cadrului-cheie curent, indexul
\texttt{kf\_index} avansează la următorul cadru-cheie și
\texttt{hold\_counter} se resetează. Dezactivarea intrării \texttt{enable}
(prin comutatorul \texttt{sw0}) resetează atât indexul cât și numărătorul,
astfel încât reactivarea pornește secvența de la început.

Ordinea servomotoarelor în vectorul de 64~biți este: degetul mic, degetul
inelar, degetul mijlociu, degetul arătător, mișcarea palmară a degetului
mare, tensiunea degetului mare, rotația încheieturii și extensia cotului.
Convențional, valoarea 0 corespunde poziției complet închise (pentru degete)
sau complet flexate (pentru cot), iar 180 corespunde poziției complet
deschise, respectiv extinse.

%==========================================================================
\section{Modulul \texttt{spi\_master} --- Interfața SPI}
\label{sec:spi-master}

Modulul \texttt{spi\_master} implementează un master SPI generic care
operează în modul~3 (CPOL=1, CPHA=1), conform cerințelor ADC-ului AD7124-8.

\subsection{Parametrii SPI mod 3}

În modul~3:
\begin{itemize}
	\item SCLK este inactiv la nivel logic '1' (CPOL=1).
	\item Datele sunt modificate pe frontul descrescător al SCLK
	      (\textit{falling edge}) și eșantionate pe frontul crescător
	      (\textit{rising edge}) --- CPHA=1.
	\item Frecvența SCLK rezultă din divizarea ceasului de sistem:
	      100\,MHz / 20 = 5\,MHz. Constanta \texttt{CLK\_DIV = 9}
	      generează o semiperioadă de 10~cicluri (numărare de la 0 la 9),
	      rezultând o perioadă completă de 20~cicluri de ceas.
\end{itemize}

\subsection{Mașina de stări}

\begin{enumerate}
	\item \textbf{S\_IDLE} --- SCLK este menținut la '1'. La activarea
	      \texttt{start}, octetul de transmis este salvat în
	      \texttt{shift\_out}, iar bitul MSB este plasat pe linia MOSI.
	      Semnalul \texttt{cs\_n} este activat la '0'.

	\item \textbf{S\_LEADING\_EDGE} --- Numărătorul de divizare
	      \texttt{div\_cnt} avansează de la 0 la \texttt{CLK\_DIV} (9).
	      La atingerea pragului, SCLK comută la '0' (front descrescător),
	      moment în care receptorul eșantionează MOSI. Se trece la
	      S\_TRAILING\_EDGE.

	\item \textbf{S\_TRAILING\_EDGE} --- Numărătorul avansează din nou
	      până la 9. La atingerea pragului, SCLK comută la '1' (front
	      crescător), moment în care masterul eșantionează MISO. Bitul
	      recepționat este introdus în \texttt{shift\_in}. Dacă toți
	      8~biții au fost transferați (\texttt{bit\_cnt = 7}), se trece
	      la S\_DONE; altfel, se avansează \texttt{bit\_cnt} și se
	      actualizează MOSI cu următorul bit.

	\item \textbf{S\_DONE} --- Octetul recepționat este copiat pe
	      \texttt{rx\_data}, semnalul \texttt{done} pulsează, iar mașina
	      revine la S\_IDLE.
\end{enumerate}

\subsection{Protecția contra metastabilității pe MISO}

Similar receptorului UART, linia MISO provine dintr-un circuit extern
(ADC-ul AD7124-8) și este asincronă față de ceasul FPGA. Două flip-flop-uri
în cascadă (\texttt{miso\_s1}, \texttt{miso\_s2}) sincronizează semnalul
înainte de eșantionare:

\begin{lstlisting}[caption={Sincronizator pe linia MISO},label={lst:miso-sync}]
process(clk)
begin
    if rising_edge(clk) then
        miso_s1 <= miso;
        miso_s2 <= miso_s1;
    end if;
end process;
\end{lstlisting}

Ordinea de transfer este MSB first (bitul cel mai semnificativ este trimis
primul), conform convenției AD7124-8.

%==========================================================================
\section{Modulul \texttt{ad7124\_ctrl} --- Controlerul ADC}
\label{sec:ad7124-ctrl}

Modulul \texttt{ad7124\_ctrl} controlează ADC-ul AD7124-8 pe 24~de biți
prin intermediul modulului \texttt{spi\_master}. Funcționarea se desfășoară
în două faze: inițializarea (o singură dată după reset) și citirea ciclică
continuă.

\subsection{Faza de inițializare}

Secvența de inițializare cuprinde următorii pași:

\begin{enumerate}
	\item \textbf{Reset hardware:} Se transmit 8~octeți de 0xFF prin SPI.
	      Conform documentației AD7124-8, transmiterea a cel puțin 64~de biți
	      consecutivi de '1' resetează registrele interne ale ADC-ului la
	      valorile implicite.

	\item \textbf{Așteptare post-reset:} Un interval de 4\,ms
	      (400\,000~cicluri de ceas) permite ADC-ului să finalizeze procedura
	      internă de resetare. Această valoare depășește specificația minimă
	      din fișa de catalog, asigurând o marjă confortabilă.

	\item \textbf{Configurare registre:} Se scriu secvențial 11~registre
	      de configurare, fiecare consistând dintr-un octet de comandă
	      (adresa registrului) urmat de 2 sau 3~octeți de date:
\end{enumerate}

\begin{table}[H]
\centering
\caption{Secvența de configurare a registrelor AD7124-8}
\label{tab:adc-config}
\begin{tabular}{@{}llp{6.5cm}@{}}
\toprule
\textbf{Registru} & \textbf{Valoare} & \textbf{Semnificație} \\
\midrule
ADC\_Control (0x01) & 0x0DC0 & Full power, conversie continuă, DATA\_STATUS activat, referință internă \\
Config\_0 (0x19) & 0x0070 & Unipolar, buffere intrare, referință internă, PGA=1 \\
Filter\_0 (0x21) & 0x060180 & Filtru sinc3, FS=384, ODR=50\,Hz \\
Channel 0 (0x09) & 0x8011 & Activat, Setup 0, AINP=0, AINM=AVSS \\
Channel 1 (0x0A) & 0x8031 & Activat, Setup 0, AINP=1, AINM=AVSS \\
Channel 2 (0x0B) & 0x8051 & Activat, Setup 0, AINP=2, AINM=AVSS \\
Channel 3 (0x0C) & 0x8071 & Activat, Setup 0, AINP=3, AINM=AVSS \\
Channel 4 (0x0D) & 0x8091 & Activat, Setup 0, AINP=4, AINM=AVSS \\
Channel 5 (0x0E) & 0x80B1 & Activat, Setup 0, AINP=5, AINM=AVSS \\
Channel 6 (0x0F) & 0x80D1 & Activat, Setup 0, AINP=6, AINM=AVSS \\
Channel 7 (0x10) & 0x80F1 & Activat, Setup 0, AINP=7, AINM=AVSS \\
\bottomrule
\end{tabular}
\end{table}

Rata de ieșire a datelor (ODR) de 50\,Hz corespunde exact ciclului de servo
de 20\,ms, asigurând că o nouă măsurătoare este disponibilă la fiecare
actualizare a servomotoarelor.

Toate cele 8~canale sunt configurate în mod \textit{single-ended}, cu
intrarea negativă conectată la AVSS (masa analogică). Fiecare canal citește
semnalul de feedback analog de la potențiometrul intern al servomotorului
MG996R corespunzător.

\subsection{Faza de citire ciclică}

După finalizarea inițializării, modulul intră într-un ciclu continuu de
citire a registrului de date. Secvența de citire este:

\begin{enumerate}
	\item Se trimite comanda de citire a registrului de date: 0x42
	      (RD=1 | adresa registrului de date = 0x02).
	\item Se recepționează 3~octeți de date (24~biți, valoarea ADC brută)
	      și 1~octet de status (care conține, în biții 3:0, numărul
	      canalului de pe care provine conversiunea).
	\item Se extrage numărul canalului din octetul de status și se efectuează
	      conversia în unghi:
\end{enumerate}

\begin{equation}
\text{angle} = \frac{\text{adc\_raw} \times 180}{2^{24}}
\label{eq:adc-to-angle}
\end{equation}

Implementarea în VHDL utilizează înmulțirea cu 180 urmată de deplasare la
dreapta cu 24~de biți (echivalent cu împărțirea la $2^{24}$). Aceasta
mapează domeniul complet al ADC-ului (0x000000 -- 0xFFFFFF) la intervalul
0°--180°.

Un registru de 8~biți \texttt{channels\_read} urmărește care canale au
fost citite. Când toate cele 8~bituri sunt setate (0xFF), modulul emite
semnalul \texttt{meas\_valid} și furnizează toți cei 8~unghiuri măsurați
pe portul \texttt{meas\_angles}, apoi resetează registrul pentru un nou
ciclu de citire.

Semnalul \texttt{config\_done} este conectat la LED-ul 0 de pe placa Basys3,
oferind o indicație vizuală că inițializarea ADC-ului s-a finalizat cu succes.

%==========================================================================
\section{Modulul \texttt{kalman\_engine} --- Filtrul Kalman}
\label{sec:kalman-engine}

Modulul \texttt{kalman\_engine} implementează nucleul algoritmului de filtrare
Kalman cu 2~stări (poziție și viteză) în aritmetică virgulă fixă Q16.16,
multiplexat în timp pe cele 8~canale de servo. Acesta este modulul central al
proiectului, responsabil pentru stabilizarea mișcărilor mâinii robotice prin
estimarea optimă a traiectoriei pe baza măsurătorilor zgomotoase.

\subsection{Constantele de configurare}

Tabelul~\ref{tab:kalman-constants} prezintă constantele utilizate de filtru,
cu valorile în reprezentare Q16.16 și echivalentele reale.

\begin{table}[H]
\centering
\caption{Constantele filtrului Kalman în reprezentare Q16.16}
\label{tab:kalman-constants}
\begin{tabular}{@{}lrll@{}}
\toprule
\textbf{Constantă} & \textbf{Valoare Q16.16} & \textbf{Valoare reală} & \textbf{Semnificație} \\
\midrule
ONE\_Q16 & 65\,536   & 1,0          & Reprezentarea unității \\
DT\_Q16  & 1\,311    & 0,02\,s      & Perioada de actualizare (20\,ms) \\
Q00      & 655       & 0,01\,grad$^2$ & Zgomot de proces --- poziție \\
Q11      & 6\,554    & 0,1\,(grad/s)$^2$ & Zgomot de proces --- viteză \\
R\_NOISE & 65\,536   & 1,0\,grad$^2$ & Zgomot de măsurare \\
P\_INIT  & 655\,360  & 10,0         & Incertitudine inițială \\
\bottomrule
\end{tabular}
\end{table}

Valorile Q00 și Q11 definesc cât de mult \textit{distrust-ăm} predicția
(valori mai mari = predicție mai nesigură, filtru mai receptiv). R\_NOISE
definește cât de mult \textit{distrust-ăm} măsurătoarea ADC (valori mai
mari = măsurătoare mai zgomotoasă, filtru mai lin). Aceste constante
pot fi ajustate în funcție de comportamentul observat pe hardware-ul real.

\subsection{Memoria de stare (State RAM)}

Starea filtrului Kalman pentru cele 8~canale este stocată într-o memorie RAM
internă de $8 \times 5 = 40$ cuvinte de câte 32~de biți:

\begin{table}[H]
\centering
\caption{Structura memoriei de stare pentru un canal}
\label{tab:state-ram}
\begin{tabular}{@{}clp{8cm}@{}}
\toprule
\textbf{Offset} & \textbf{Variabilă} & \textbf{Descriere} \\
\midrule
+0 & $x_0$ & Poziția estimată (grade, Q16.16) \\
+1 & $x_1$ & Viteza estimată (grade/secundă, Q16.16) \\
+2 & $P_{00}$ & Incertitudinea poziției \\
+3 & $P_{01}$ & Corelația poziție-viteză \\
+4 & $P_{11}$ & Incertitudinea vitezei \\
\bottomrule
\end{tabular}
\end{table}

Adresa de bază pentru canalul $k$ este $5k$, iar funcția \texttt{ch\_base}
calculează această adresă. La inițializare, $x_0$ și $x_1$ sunt setate la
zero, iar $P_{00}$ și $P_{11}$ sunt setate la P\_INIT = 10,0 (incertitudine
inițială mare, semnificând că filtrul nu are încă informații despre poziția
reală a servomotoarelor).

\subsection{Funcția de multiplicare Q16.16}

Multiplicarea a două numere Q16.16 produce un rezultat pe 64~de biți cu
32~de biți fracționari. Pentru a reveni la formatul Q16.16 (16~biți
fracționari), se extrag biții 47:16 din produsul complet:

\begin{lstlisting}[caption={Funcția de multiplicare în virgulă fixă Q16.16},label={lst:mul-q16}]
function mul_q16(a : signed(31 downto 0);
                 b : signed(31 downto 0))
    return signed is
    variable product : signed(63 downto 0);
begin
    product := a * b;
    return product(47 downto 16);
end function;
\end{lstlisting}

Această operație este echivalentă matematic cu $(a \times b) \gg 16$,
menținând precizia fracționară fără a necesita unitate de virgulă mobilă.
Sintetizatorul Vivado mapează aceste multiplicări pe blocurile DSP48E1
disponibile în Artix-7, eliminând necesitatea implementării cu logică
combinațională.

\subsection{Mașina de stări a filtrului Kalman}
\label{subsec:kalman-fsm}

Filtrul procesează secvențial cele 8~canale de servo, executând pentru
fiecare canal un ciclu complet de predicție și actualizare. Mașina de
stări cuprinde 15~stări, grupate funcțional:

\subsubsection{Inițializare (S\_INIT\_RAM)}

La pornire sau după reset, se parcurge memoria de stare de la adresa 0 la
39, inițializând fiecare cuvânt conform structurii din
Tabelul~\ref{tab:state-ram}. Pozițiile de offset 2 și 4 (P$_{00}$ și
P$_{11}$) primesc valoarea P\_INIT, iar celelalte sunt inițializate la zero.
După finalizare, semnalul \texttt{initialized} este activat și mașina trece
la S\_IDLE.

\subsubsection{Încărcare (S\_LOAD)}

Se citesc cele 5~variabile de stare ale canalului curent din RAM și se
convertește unghiul măsurat din format pe 8~biți în Q16.16 (multiplicare
cu 65\,536). Funcția \texttt{get\_angle\_q16} extrage octetul corespunzător
canalului curent din vectorul de 64~biți al măsurătorilor.

\subsubsection{Predicție (S\_PREDICT\_X, S\_PREDICT\_P00, P01, P11)}

Se aplică ecuațiile de predicție ale filtrului Kalman:

\begin{align}
x_{0,\text{pred}} &= x_0 + dt \cdot x_1 \label{eq:predict-x0} \\
x_{1,\text{pred}} &= x_1 \label{eq:predict-x1}
\end{align}

Viteza nu se modifică în predicție deoarece modelul cu 2~stări nu include
accelerația. Predicția covarinaței crește incertitudinea:

\begin{align}
P_{00}' &= P_{00} + 2 \cdot dt \cdot P_{01} + dt^2 \cdot P_{11} + Q_{00}
\label{eq:predict-p00} \\
P_{01}' &= P_{01} + dt \cdot P_{11} \label{eq:predict-p01} \\
P_{11}' &= P_{11} + Q_{11} \label{eq:predict-p11}
\end{align}

Implementarea VHDL a predicției P$_{00}$ utilizează funcția \texttt{mul\_q16}
pentru toate produsele:

\begin{lstlisting}[caption={Predicția covarianței P00 în VHDL},label={lst:predict-p00}]
when S_PREDICT_P00 =>
    pp00 <= p00 + mul_q16(DT_Q16, p01) +
            mul_q16(DT_Q16, p01) +
            mul_q16(DT_Q16, mul_q16(DT_Q16, p11)) + Q00;
    fsm <= S_PREDICT_P01;
\end{lstlisting}

\subsubsection{Actualizare (S\_UPDATE\_S, K0, K1, X, P)}

Inovația (\textit{innovation}) reprezintă diferența dintre măsurătoarea
reală și predicție:

\begin{equation}
y = z_{\text{meas}} - x_{0,\text{pred}}
\label{eq:innovation}
\end{equation}

Covarianța inovației combină incertitudinea predicției cu zgomotul de
măsurare:

\begin{equation}
S = P_{00}' + R
\label{eq:innovation-cov}
\end{equation}

Câștigul Kalman determină ponderea acordată măsurătorii vs. predicției:

\begin{align}
K_0 &= P_{00}' / S \label{eq:kalman-k0} \\
K_1 &= P_{01}' / S \label{eq:kalman-k1}
\end{align}

Dacă incertitudinea predicției ($P_{00}'$) este mare comparativ cu zgomotul
de măsurare ($R$), câștigul se apropie de 1,0 --- filtrul încredințează
mai mult măsurătorii. Dacă $P_{00}'$ este mic, câștigul se apropie de 0 ---
filtrul încredințează mai mult predicția.

Implementarea diviziunii în VHDL utilizează operatorul \texttt{/}
aplicat pe valori extinse la 64~de biți, cu deplasare la stânga prealabilă
a numărătorului cu 16~biți pentru a menține precizia Q16.16:

\begin{lstlisting}[caption={Calculul câștigului Kalman K0 cu protecție împotriva împărțirii la zero},label={lst:kalman-k0}]
when S_UPDATE_K0 =>
    if s_val /= 0 then
        tmp64 := shift_left(resize(pp00, 64), 16);
        k0 <= resize(tmp64 / resize(s_val, 64), 32);
    else
        k0 <= ONE_Q16;
    end if;
\end{lstlisting}

Protecția împotriva împărțirii la zero setează $K_0 = 1,0$ în cazul extrem
în care $S = 0$, ceea ce ar însemna încredere totală în măsurătoare.

Corecția stării:

\begin{align}
x_0 &= x_{0,\text{pred}} + K_0 \cdot y \label{eq:update-x0} \\
x_1 &= x_{1,\text{pred}} + K_1 \cdot y \label{eq:update-x1}
\end{align}

Ecuația~\eqref{eq:update-x1} este mecanismul prin care viteza \textit{învață}:
dacă măsurătoarea arată constant că servomotorul s-a deplasat mai mult decât
a prezis predicția, inovația $y$ este pozitivă, iar $K_1 \cdot y$ crește
estimarea vitezei. Reciproc, dacă servomotorul este în repaus, inovația
scade, iar viteza converge spre zero.

Actualizarea covarianței (incertitudinea scade după integrarea măsurătorii):

\begin{align}
P_{00} &= (1 - K_0) \cdot P_{00}' \label{eq:update-p00} \\
P_{01} &= (1 - K_0) \cdot P_{01}' \label{eq:update-p01} \\
P_{11} &= -K_1 \cdot P_{01}' + P_{11}' \label{eq:update-p11}
\end{align}

\subsubsection{Stocare și avansare (S\_STORE, S\_NEXT\_CH, S\_OUTPUT)}

Cele 5~variabile de stare actualizate sunt scrise înapoi în RAM la adresa
de bază a canalului curent. Poziția estimată ($x_0$) este convertită de la
Q16.16 la un unghi pe 8~biți prin deplasare la dreapta cu 16~biți, urmată
de limitare la intervalul 0--180:

\begin{lstlisting}[caption={Conversia Q16.16 la unghi pe 8 biți cu limitare},label={lst:angle-clamp}]
result_angle := to_integer(shift_right(x0, 16));
if result_angle < 0 then
    result_angle := 0;
elsif result_angle > 180 then
    result_angle := 180;
end if;
\end{lstlisting}

După procesarea tuturor celor 8~canale ($\texttt{channel} = 0, 1, \ldots, 7$),
starea S\_OUTPUT transferă toți unghiurile filtrate pe portul
\texttt{output\_angles} și activează semnalul \texttt{output\_valid}
pentru un ciclu de ceas. Mașina revine apoi la S\_IDLE, așteptând
următorul tick de 20\,ms.

\subsection{Buget temporal}

Fiecare canal necesită aproximativ 14~multiplicări, 2~diviziuni și 14~adunări,
totalizând circa 182~cicluri de ceas. Procesarea tuturor celor 8~canale
necesită aproximativ $8 \times 182 = 1\,456$ cicluri din cele 2\,000\,000
disponibile între două tick-uri de 20\,ms, reprezentând un grad de utilizare
de doar 0,07\%. Aceasta lasă un buget de ceas generos pentru eventuale
extinderi viitoare, cum ar fi trecerea la un filtru Kalman cu 3~stări.

%==========================================================================
\section{Modulul \texttt{pwm\_gen} --- Generatorul PWM}
\label{sec:pwm-gen}

Modulul \texttt{pwm\_gen} generează un semnal PWM standard pentru un
servomotor. Standardul servo definește:

\begin{itemize}
	\item Perioadă: 20\,ms (50\,Hz)
	\item Puls de 1,0\,ms $\rightarrow$ 0° (poziție minimă)
	\item Puls de 2,0\,ms $\rightarrow$ 180° (poziție maximă)
\end{itemize}

La frecvența de ceas de 100\,MHz:
\begin{itemize}
	\item 1,0\,ms = 100\,000 de cicluri
	\item 2,0\,ms = 200\,000 de cicluri
	\item Intervalul util = 100\,000 de cicluri pentru 180°
	\item Factor de conversie: $100\,000 / 180 \approx 556$
\end{itemize}

Pragul PWM se calculează ca:

\begin{equation}
\text{threshold} = 100\,000 + \text{angle} \times 556
\label{eq:pwm-threshold}
\end{equation}

Ieșirea PWM este '1' cât timp numărătorul intern (pe 21~de biți, numărând de
la 0 la 1\,999\,999) este mai mic decât pragul, și '0' în rest. Acest
principiu de comparație este implementat printr-o singură instrucțiune
combinațională:

\begin{lstlisting}[caption={Generarea semnalului PWM prin comparație},label={lst:pwm-compare}]
pwm_out <= '1' when counter < threshold else '0';
\end{lstlisting}

\subsection{Generarea a 8 instanțe}

În modulul de nivel superior \texttt{basys3\_top}, se generează 8~instanțe
ale modulului \texttt{pwm\_gen} prin intermediul instrucțiunii
\texttt{for...generate}:

\begin{lstlisting}[caption={Generarea a 8 instanțe PWM prin instrucțiunea generate},label={lst:gen-pwm}]
gen_pwm : for i in 0 to 7 generate
    u_pwm : entity work.pwm_gen
        port map (
            clk     => clk,
            rst     => rst,
            angle   => kalman_out((7-i)*8+7 downto (7-i)*8),
            pwm_out => pwm_out(i)
        );
end generate;
\end{lstlisting}

Fiecare instanță primește octetul de unghi corespunzător din vectorul de
64~biți al ieșirii filtrului Kalman. Indexarea $(7-i) \times 8$ asigură
corespondența corectă între canalele logice (0--7) și pozițiile din vectorul
de biți.

Deoarece fiecare instanță PWM posedă propriul numărător hardware, toate
cele 8~canale funcționează în paralel, fără multiplexare în timp. Această
proprietate a implementării FPGA garantează că actualizarea unui canal
nu afectează în niciun fel celelalte canale, eliminând jitter-ul de
interleaving prezent în implementările bazate pe microcontroler.

%==========================================================================
\section{Modulul \texttt{basys3\_top} --- Integrarea la nivel superior}
\label{sec:basys3-top}

Modulul \texttt{basys3\_top} este entitatea de nivel superior care
instanțiază și interconectează toate cele 12~module descrise anterior.
Porturile sale corespund pinilor fizici ai plăcii Basys3 Artix-7.

\subsection{Conectarea porturilor}

\begin{table}[H]
\centering
\caption{Maparea principală a porturilor la funcții}
\label{tab:port-mapping}
\begin{tabular}{@{}llp{6cm}@{}}
\toprule
\textbf{Port} & \textbf{Pin Basys3} & \textbf{Funcție} \\
\midrule
clk   & W5  & Oscilator 100\,MHz \\
sw0   & V17 & Comutator mod demo \\
btnC  & U18 & Buton reset (centru) \\
RsRx  & B18 & UART recepție \\
RsTx  & A18 & UART transmisie \\
JA0--JA3 & J1, L2, J2, G2 & PWM servouri 0--3 (pinky, ring, middle, index) \\
JB0--JB3 & A14, A16, B15, B16 & PWM servouri 4--7 (thumb, wrist, elbow) \\
JC0 (CS) & K17 & SPI chip select AD7124 \\
JC1 (MOSI) & M18 & SPI date ieșire \\
JC2 (MISO) & N17 & SPI date intrare \\
JC3 (SCLK) & P18 & SPI ceas \\
led[0:15] & diverse & Indicatoare de stare \\
\bottomrule
\end{tabular}
\end{table}

\subsection{Gestionarea pinilor neutilizați}

Toți pinii neutilizați ai plăcii Basys3 sunt configurați la stări sigure
pentru a preveni intrări flotante și consum inutil de curent:

\begin{itemize}
	\item \textbf{JA4--JA7, JB4--JB7, JC4--JC7:} Jumătatea superioară a
	      conectorilor Pmod, forțată la '0' (drive low).
	\item \textbf{seg[6:0], dp, an[3:0]:} Afișajul cu 7~segmente este activ
	      la nivel logic scăzut; se forțează la '1' pentru a dezactiva toate
	      segmentele și cifrele.
	\item \textbf{vgaRed, vgaGreen, vgaBlue, Hsync, Vsync:} Ieșirile VGA
	      sunt forțate la '0'.
	\item \textbf{PS2Clk, PS2Data:} Interfața PS/2 este lăsată în
	      înaltă impedanță ('Z') deoarece este un bus open-drain.
	\item \textbf{JXADC[7:0]:} Header-ul analogic este lăsat în
	      înaltă impedanță ('Z').
	\item \textbf{sw[15:1], btnU/D/L/R:} Comutatoarele și butoanele
	      neutilizate rămân ca intrări, citite dar neasignate.
\end{itemize}

\subsection{Indicatoarele LED}

Cele 16~LED-uri de pe placa Basys3 sunt utilizate ca indicatoare de stare:

\begin{table}[H]
\centering
\caption{Asignarea indicatoarelor LED}
\label{tab:led-assignment}
\begin{tabular}{@{}cl@{}}
\toprule
\textbf{LED} & \textbf{Semnal indicat} \\
\midrule
0 & Configurare ADC finalizată (\texttt{adc\_cfg\_done}) \\
1 & Măsurătoare ADC validă (\texttt{meas\_valid}) \\
2 & Cadru UART recepționat (\texttt{frame\_valid}) \\
3 & Eroare de checksum (\texttt{frame\_chksum\_err}) \\
4 & Filtru Kalman în execuție (\texttt{kalman\_busy}) \\
5 & Mod demo activ, reflectă sw0 \\
6 & Transmițător UART ocupat (\texttt{uart\_tx\_busy}) \\
7 & SPI ocupat (\texttt{spi\_busy\_sig}) \\
8--15 & Unghiul filtrat al servo 0 (8 biți) \\
\bottomrule
\end{tabular}
\end{table}

LED-urile 8--15 afișează în binar unghiul filtrat al primului servo
(degetul mic), oferind o verificare vizuală imediată a funcționării
filtrului Kalman.

\subsection{Logica de declanșare a răspunsului}

Transmisia unui cadru de răspuns (readback) este declanșată atunci când
un cadru valid cu direcția 0xFE este recepționat:

\begin{lstlisting}[caption={Logica de declanșare a cadrului de răspuns},label={lst:readback-trigger}]
process(clk)
begin
    if rising_edge(clk) then
        send_readback <= '0';
        if frame_valid = '1' and frame_dir = x"FE" then
            send_readback <= '1';
        end if;
    end if;
end process;
\end{lstlisting}

Această logică sincronă generează un puls de un ciclu de ceas pe semnalul
\texttt{send\_readback}, care declanșează modulul \texttt{frame\_tx} să
transmită un cadru de răspuns conținând cele 8~unghiuri filtrate curente.

%==========================================================================
\section{Proiectarea hardware --- circuitele de interfață}
\label{sec:hw-design}

Pe lângă implementarea FPGA, sistemul necesită circuite auxiliare de
interfață pentru izolarea galvanică, alimentarea servomotoarelor și
conversia de nivel logic.

\subsection{Izolarea galvanică a semnalelor PWM --- ADUM1400}

Semnalele PWM generate de FPGA (la nivel logic de 3,3\,V) trebuie
convertite la 5\,V pentru servomotoarele MG996R. Se utilizează două circuite
integrate ADUM1400, fiecare oferind 4~canale de izolare galvanică
digitală. ADUM1400 este un izolator bazat pe tehnologia \textit{iCoupler}
(transformatoare microfabricate), oferind:

\begin{itemize}
	\item Izolare galvanică de 2\,500\,V\textsubscript{rms}
	\item Timp de propagare tipic de 32\,ns (neglijabil pentru semnale PWM)
	\item Conversie de nivel logic intrinsecă: alimentat la 3,3\,V pe partea
	      FPGA și la 5\,V pe partea servomotoare
\end{itemize}

Primul ADUM1400 convertește semnalele PWM pentru servomotoarele 0--3
(JA0--JA3), iar al doilea pentru servomotoarele 4--7 (JB0--JB3).

\subsection{Alimentarea servomotoarelor --- DC2468A}

Cele 8~servomotoare MG996R consumă fiecare până la 700\,mA sub sarcină,
totalizând un consum potențial de 5,6\,A la 5\,V. Sursa de alimentare
utilizată este modulul DC2468A (bazat pe regulatorul LT3045), care
convertește 12\,V dintr-un alimentator extern la 5\,V cu un curent maxim
de 8\,A. Această marjă de 30\% asigură funcționarea stabilă chiar și
când toate servomotoarele sunt acționate simultan.

\subsection{Alimentarea circuitului ADC --- ADP165}

ADC-ul AD7124-8 și izolorul SPI MAX14850 necesită alimentare la 3,3\,V
separată de domeniul FPGA. Se utilizează un regulator LDO ADP165 care
furnizează 3,3\,V cu zgomot ultra-scăzut din sursa de 5\,V, minimizând
impactul zgomotului de alimentare asupra preciziei măsurătorilor ADC.

\subsection{Izolarea SPI --- MAX14850}

Interfața SPI dintre FPGA și AD7124-8 traversează un izolator digital
MAX14850 cu 4~canale. Acesta asigură izolarea galvanică între domeniul
de ceas FPGA (3,3\,V) și domeniul ADC (3,3\,V alimentat separat),
prevenind buclele de masă care ar introduce zgomot în măsurătorile ADC.
Cele 4~canale sunt utilizate pentru: SCLK, MOSI, MISO și CS.

\subsection{Conectarea servomotoarelor MG996R}

Servomotoarele MG996R sunt conectate prin cabluri cu 3~fire:
\begin{itemize}
	\item \textbf{Maro:} GND (masa comună)
	\item \textbf{Roșu:} V+ (5\,V de la DC2468A)
	\item \textbf{Portocaliu:} Semnal PWM (5\,V de la ADUM1400)
\end{itemize}

Fiecare MG996R dispune și de un al patrulea fir (alb), conectat la
potențiometrul intern de feedback analogic. Acest semnal este direcționat
către una din cele 8~intrări analogice ale AD7124-8 pentru citirea
poziției reale a servomotorului.
